\documentclass{article}
\usepackage[indonesian]{babel}
\begin{document}
Galat baku penduga rasio dihitung dengan
\begin{equation}
\mathrm{SE}(\hat R)=\sqrt{\frac{1-f}{n}\,s_r^2},
\end{equation}

dengan $n$ menyatakan ukuran sampel dan $f$ fraksi penarikan sampel.

Ragam sampel diduga dengan

\[ s^2=\frac{1}{n-1}\sum_{i=1}^{n}(x_i-\bar x)^2 \]
Dengan $\bar x$ menyatakan rata-rata sampel.

Model regresi linear ditulis sebagai
\begin{equation}
y_i=\beta_0+\beta_1x_i+\varepsilon_i,
\end{equation}
dengan $\varepsilon_i$ galat acak yang saling bebas.
Sifat penduganya dibahas pada subbab berikut.
\end{document}
